\begin{definition}[A native endpoint route beside the partner plaquette]
    \label{def:peps_torus_swept_physical_endpoint_route}
    \lean{TNLean.PEPS.sweptPhysicalBRouteCoordinate,
        TNLean.PEPS.torusSweptPhysicalBRouteVertex,
        TNLean.PEPS.torusSweptPhysicalBRoute}
    \leanok
    \uses{def:peps_torus_swept_patch_endpoint_geometry,
        def:peps_torus_swept_patch_endpoint_walks}
    Let the horizontal and vertical torus periods be at least eight and seven,
    respectively. Fix any translated position $v$ and write
    $p(i,j)=(v_x+i,v_y+j)$. Consider the successive plaquette bases
    \begin{align}
         & p(1,1),p(2,1),p(2,2),p(3,2),p(4,2),p(4,1),p(4,0), \\
         & p(3,0),p(2,0),p(1,0),p(0,0),p(0,-1),p(0,-2).
    \end{align}
    Each successive pair is joined by an actual unit bond of the native
    torus graph. Their ordered traversal defines a twelve-step walk $\rho_v$.
    This chosen route passes above, to the right of, and below the $A_1$
    plaquette before continuing towards $B_2$.
    It is an auxiliary finite geometry for the physical movement discussed in
    \cite[fluxon-braiding passage]{Schuch2010PEPS}, lines 2340--2423.
\end{definition}

\begin{theorem}[Actual endpoints, partner avoidance, and final adjacency]
    \label{thm:peps_torus_swept_physical_endpoint_route}
    \lean{TNLean.PEPS.torusSweptPhysicalBRoute_length,
        TNLean.PEPS.torusSweptPhysicalBRoute_start,
        TNLean.PEPS.torusSweptPhysicalBRoute_end,
        TNLean.PEPS.torusSweptPhysicalBRoute_end_adj_partner,
        TNLean.PEPS.torusSweptPhysicalBRoute_avoids_partners}
    \leanok
    \uses{def:peps_torus_swept_physical_endpoint_route,
        thm:peps_torus_swept_patch_endpoint_geometry}
    The actual walk $\rho_v$ has length twelve. It starts at
    $B_1=p(1,1)$ and ends at $p(0,-2)$, immediately left of
    $B_2=p(1,-2)$. Its final base and $B_2$ are adjacent in the native graph.
    None of its thirteen bases equals $A_2=p(-1,1)$, $A_1=p(3,1)$, or $B_2$.
    These assertions hold at every translated position, including either
    periodic seam; no adjacency or endpoint-avoidance premise is supplied.
\end{theorem}
\begin{proof}\leanok
    Evaluate the first and last coordinates and count the twelve successive
    unit steps. The final adjacency is the native rightward bond. Distinct
    offsets in the coordinate envelope remain distinct modulo the periods,
    and the finite route list contains none of the three fixed partner bases.
\end{proof}

\begin{definition}[The thirty-site movement footprint]
    \label{def:peps_torus_swept_physical_route_footprint}
    \lean{TNLean.PEPS.sweptPhysicalRouteFootprintCoordinates,
        TNLean.PEPS.torusSweptPhysicalRouteCoordinate,
        TNLean.PEPS.torusSweptPhysicalRouteFootprint}
    \leanok
    \uses{def:peps_torus_swept_physical_endpoint_route,
        def:peps_torus_swept_patch_endpoint_geometry}
    Let $F_v$ be the union of the twenty-site four-endpoint witness $R_v$
    and the four physical corners of each of the thirteen route plaquettes.
    All its vertices have offsets
    $-1\leq i\leq5$ and $-2\leq j\leq3$ from $v$.
    This seven-column, six-row envelope contains both plaquettes of every
    successive elementary movement step.
\end{definition}

\begin{theorem}[Thirty distinct sites contain the full actual route]
    \label{thm:peps_torus_swept_physical_route_footprint}
    \lean{TNLean.PEPS.card_sweptPhysicalRouteFootprintCoordinates,
        TNLean.PEPS.torusSweptPhysicalRouteCoordinate_injective,
        TNLean.PEPS.torusSweptPhysicalRouteFootprint_card,
        TNLean.PEPS.torusSweptPhysicalBRoute_corner_mem_footprint,
        TNLean.PEPS.torusSweptPhysicalBRoute_vertex_mem_footprint,
        TNLean.PEPS.torusSweptPhysicalBRoute_mem_footprint}
    \leanok
    \uses{def:peps_torus_swept_physical_route_footprint,
        thm:peps_torus_swept_patch_endpoint_geometry}
    The integer-coordinate footprint has thirty sites. Its translated native
    realization is injective on the entire seven-column, six-row envelope,
    and hence $|F_v|=30$. Every corner of every route plaquette belongs to
    $F_v$, as does every vertex of the actual walk $\rho_v$.
    Thus the elementary movement footprints fit into one fixed finite region.

    This result supplies geometry only. It does not assert the composition
    of the physical movement operators, the resulting bond assignment,
    a winding-number calculation, or the prescribed physical braid.
\end{theorem}
\begin{proof}\leanok
    Count the union of the finite corner lists. Each offset range is shorter
    than its period, so equality of residue classes implies equality of
    coordinates, including after translation across either seam.
    Each route corner belongs to the defining corner union; the route vertices
    are its lower-left corners. The explicit walk support is the route list.
\end{proof}
